\section{Appendix C — Core principles}

\subsection{Be consistent}

The dungeon is unpredictable, but its underlying rules should not be. Players must be able to learn from experience.

\subsection{Reward observation}

Descriptions, door appearances, shifting sounds and notes are valuable sources of information.

\subsection{Make retreat viable}

Not every encounter is intended to be defeated. Escaping a powerful enemy is a legitimate tactical victory.

\subsection{Preserve player agency}

Dangerous traps and encounters should offer meaningful opportunities to react, escape or experiment.

\subsection{Keep rewards meaningful}

Magical objects should create strategic possibilities rather than merely increase combat statistics.

\subsection{Let the dungeon create the difficulty}

Shifting geography, resetting encounters, limited resources and permanent collapse provide pressure without arbitrary intervention.

\subsection{Distinguish danger from unfairness}

Characters may die because of decisions, but should rarely die because the DM failed to communicate something they could reasonably have known.

Ultimately, Shifting Dungeon is not about defeating everything inside it. It is about understanding a hostile environment well enough to survive it.

{\small\itshape Design note: Structural mechanics follow the rules established during development. Party level, encounter statistics, damage scaling, resting restrictions, resurrection activation timing, unrestricted teleportation and individual escape conditions are recommended DM rulings rather than finalized rules. They remain adaptable to the party and D\&D edition.\par}
